\section{Boundary Current Applications}\label{manual:boundary-current-applications}
\aodnoteblock{Data-support class}{Boundary-current applications are D0 exact scoped-flux fixtures unless an external report map is declared. Errors are Stokes residuals, moving-boundary identity residuals, current-split residuals, and sheddic-route conservation residuals.}
\aodnoteblock{Scope}{This section applies the same scoped AFC flux calculation to wire, fluid, plasma, shell-electromagnetic comparison, and phase-fluid SADAR records. The declared boundary selects the application case of the declared-scope calculus.}

\subsection{Shared scoped-flux setup}\label{manual:boundary-current:shared-setup}
Declare an application scope
\[
\mathbf a\vdash(F,R,\partial R,\omega),
\qquad
B=(\partial F,\omega).
\]
The scoped AFC flux is
\[
\phi_{ij}(B)=p^D_{ij}(B)A_{ij}(B),
\]
with
\[
p^D_{ij}(B)=C_{3,ij}\rho^D_{\omega,ij}(B),
\qquad
\rho^D_{\omega,ij}(B)=\min(\rho^D_{ij}(B),|\omega|_{ij}),
\]
\[
A_{ij}(B)=\Delta\chi^{\mathrm{lag}}_{ij}(B)i_{ij}(B).
\]
The slice flux and window current are
\[
\Phi^D_{\partial R}(t;B|_t)
=
\sum_{e\in\partial R(t)}\phi_e(B|_t),
\]
\[
J^D_{\partial R}(\omega;B)
=
\sum_{t\in\omega}\Phi^D_{\partial R}(t;B|_t).
\]
For moving boundaries use
\[
\Phi^D_{\partial R,\mathrm{mov}}(t)
=
\Phi^D_{\mathrm{cross}}(t)
-
\Phi^D_{\mathrm{sweep}}(t)
+
\Phi^D_{\mathrm{shedding}}(t).
\]
The directional SADAR quantity remains
\[
\mathbf A_F(B)=\sum_{(i,j)\in F(B)}\phi_{ij}(B),
\qquad
\operatorname{motion}_F(B)=\mathbf A_F(B).
\]

\begin{figure}[h]
\centering
\scriptsize
\begin{tikzpicture}[node distance=0.55cm, every node/.style={draw,rounded corners,align=center,inner sep=4pt}]
\node (scope) {declared scope\\$\mathbf a\vdash(F,R,\partial R,\omega)$};
\node[below=of scope] (flux) {scoped flux\\$\phi_{ij}(B)=p^D_{ij}(B)A_{ij}(B)$};
\node[below left=0.65cm and 1.25cm of flux] (slice) {slice flux\\$\Phi^D_{\partial R}(t;B|_t)$};
\node[below right=0.65cm and 1.25cm of flux] (moving) {moving boundary\\$\Phi^D_{\mathrm{cross}}-\Phi^D_{\mathrm{sweep}}+\Phi^D_{\mathrm{shedding}}$};
\node[below=1.0cm of flux] (window) {window current\\$J^D_{\partial R}(\omega;B)$};
\node[below=of window] (sadar) {SADAR attention\\$\mathbf A_F(B)=\sum\phi_{ij}(B)$};
\draw[-{Latex}] (scope) -- (flux);
\draw[-{Latex}] (flux) -- (slice);
\draw[-{Latex}] (flux) -- (moving);
\draw[-{Latex}] (slice) -- (window);
\draw[-{Latex}] (moving) -- (window);
\draw[-{Latex}] (window) -- (sadar);
\end{tikzpicture}
\caption{Boundary-current application spine. The same scoped flux produces slice flux, moving-boundary flux, window current, and the SADAR attention vector on the declared boundary/window.}
\label{fig:manual-boundary-current-spine}
\end{figure}

\subsection{Canonical worked substitution}\label{manual:boundary-current:worked-substitution}
Use the signed moving-boundary calculation
\[
\Phi^D_{\mathrm{cross}}=8,
\qquad
\Phi^D_{\mathrm{sweep}}=3,
\qquad
\Phi^D_{\mathrm{shedding}}=2.
\]
Then
\[
\Phi^D_{\partial R,\mathrm{mov}}
=
8-3+2=7.
\]
For a channel \((1,2)\), declare
\[
C_{3,12}=3,
\qquad
\rho^D_{12}(B)=5,
\qquad
|\omega|_{12}=4,
\qquad
A_{12}(B)=\frac13.
\]
Then
\[
\rho^D_{\omega,12}(B)=\min(5,4)=4,
\]
\[
p^D_{12}(B)=3\cdot4=12,
\]
\[
\phi_{12}(B)=12\cdot\frac13=4,
\qquad
\phi_{21}(B)=-4.
\]

\begin{table}[h]
\centering
\scriptsize
\caption{Canonical scoped-flux substitution.}
\begin{tabular}{l c}
\toprule
Quantity & Value\\
\midrule
$\Phi^D_{\mathrm{cross}}$ & $8$\\
$\Phi^D_{\mathrm{sweep}}$ & $3$\\
$\Phi^D_{\mathrm{shedding}}$ & $2$\\
$\Phi^D_{\partial R,\mathrm{mov}}$ & $8-3+2=7$\\
$\rho^D_{\omega,12}(B)$ & $\min(5,4)=4$\\
$p^D_{12}(B)$ & $12$\\
$A_{12}(B)$ & $1/3$\\
$\phi_{12}(B)$ & $4$\\
$\phi_{21}(B)$ & $-4$\\
\bottomrule
\end{tabular}
\label{tab:manual-boundary-current-substitution}
\end{table}

\subsection{Wire boundary current}\label{manual:boundary-current:wire}
Declare
\[
\mathbf a_{\mathrm{wire}}
\vdash
(F_{\mathrm{wire}},R_{\mathrm{core}},\partial R_{\mathrm{shell}},\omega_{\mathrm{wire}}),
\qquad
B_{\mathrm{wire}}=(\partial F_{\mathrm{wire}},\omega_{\mathrm{wire}}).
\]
Here \(F_{\mathrm{wire}}\) is the conductor plus retained field shell, \(R_{\mathrm{core}}\) is the conductor/core region, and \(\partial R_{\mathrm{shell}}\) is the shell boundary.
The wire record computes
\[
\Phi^D_{\partial R}(t;B|_t),
\quad
J^D_{\partial R}(\omega;B),
\quad
P^D_F(B),
\quad
\mathbf A_F(B),
\quad
X_{\mathrm{shedding}}(B).
\]
\aodremark{The wire case is a boundary-duon-current example. The surrounding field shell carries returned-current duration, pressure, and SADAR attention balance.}

\subsection{Fluid interface dynamics}\label{manual:boundary-current:fluid}
Declare
\[
\mathbf a_{\mathrm{fluid}}
\vdash
(F_{\mathrm{fluid}},R_{\mathrm{water}},\partial R_{\mathrm{interface}},\omega_{\mathrm{fluid}}),
\qquad
B_{\mathrm{fluid}}=(\partial F_{\mathrm{fluid}},\omega_{\mathrm{fluid}}).
\]
The moving interface uses
\[
\Phi^D_{\partial R,\mathrm{mov}}(t)
=
\Phi^D_{\mathrm{cross}}(t)
-
\Phi^D_{\mathrm{sweep}}(t)
+
\Phi^D_{\mathrm{shedding}}(t).
\]
Here cross is current crossing the interface, sweep is current displaced by the moving boundary, and shedding is excess crossing or re-routing through the interface.
The output record carries interface current, interface pressure, SADAR attention vector, and sheddic path.

\subsection{Plasma shell dynamics}\label{manual:boundary-current:plasma}
Declare
\[
\mathbf a_{\mathrm{plasma}}
\vdash
(F_{\mathrm{plasma}},R_{\mathrm{core}},\partial R_{\mathrm{magnetic}},\omega_{\mathrm{plasma}}),
\qquad
B_{\mathrm{plasma}}=(\partial F_{\mathrm{plasma}},\omega_{\mathrm{plasma}}).
\]
The plasma case uses the same scoped-flux calculation as the fluid interface, with the declared boundary now a shell boundary. The calculation changes by boundary declaration, not by introducing a new current object.

\subsection{Shell electromagnetic comparison}\label{manual:boundary-current:shell-em}
The shell electromagnetic record is a downstream comparison record. It computes A\(\Omega\) quantities first:
\[
J^D_{\partial R}(\omega;B),
\qquad
P^D_F(B),
\qquad
\mathbf A_F(B),
\qquad
X_{\mathrm{shedding}}(B).
\]
Only after these A\(\Omega\) quantities are declared may a measured-sector comparison be attached.

\subsection{Phase-fluid SADAR rebalancing}\label{manual:boundary-current:phase-fluid}
Phase-fluid motion is the SADAR attention vector of the declared interface field:
\[
\operatorname{motion}_F(B)=\mathbf A_F(B).
\]
The local sequence is
\[
\text{phase bias}
\longrightarrow
\mathbf A_F(B)
\longrightarrow
\text{SADAR balance}
\longrightarrow
\text{shedding / reclosure}.
\]
Push and pull are signed phase-biased components of the same vector.

\subsection{Audit}\label{manual:boundary-current:falsification}
The boundary-current application record carries the declared scope, antisymmetry \(\phi_{ji}(B)=-\phi_{ij}(B)\), the signed window current as the sum of slice fluxes, and the moving-boundary cross, sweep, and shedding terms.
\[
J^D_{\partial R}(\omega;B)
\ne
\sum_{t\in\omega}\Phi^D_{\partial R}(t;B|_t)
\]
marks a window-current audit residual on the declared scope.
\aodprovenance{main:eq:b-scoped-flux; main:eq:rcd-coupling; main:eq:window-clipped-rho; main:eq:sadar-attention-vector; main:eq:moving-boundary.}
